\section{18 августа 1914 г.}

\lp{20260604_194429422}

Мой дорогой, милый, Миля! Незная дошли ли до вас две телеграммы, одно мое
закрытое и одно открытое письмо, еще раз сообщаю вам что все здоровы. Папа наш
нашелся. Он в Берлине, но без денег, и ему теперь выслали. Маня вернулась,
проехав домой почти полмесяца. От Ильиных я сегодня получила открытку из
Одессы, куда они ехали 3 недели и очень рискованно. О Мариетте ничего не знаю.
Каждый день читаю списки русских и не встречаю ее имени. От вас получены две
телеграммы, шедшие до нас, очевидно, бесконечно долго. Папаша получил письмо,
но уже не к рождению, а на две недели позже. Я получила ваше письмо 16-го (от
8-го новым стилем). Значит оно шло 20 дней. Последнее письмо из Мюнхена было,
когда объявление войны было только между

\lp{20260604_194437600}

\noindent
Австрией и Сербией. Ильины очень беспокоятся о вас. Их адрес: Крым, Судак, дача
Черцых Ани. Поехали отдохнуть. Марго в Москве. Хлопочет по устройству у себя
лазарета. Мы все помогаем чем можем. Колю призывали, но признали негодным и
переписали из ратников 1-го ополчения в ратники 2-го ополчения. Сабуров
командует дружиной в Калуге. Зяська посещает какой-то инженерный корпус, но его
вероятно освободят, т.к. их слишком много.

Напишите мне, умоляю, о себе! Как можно подробнее! Неужели Ядвига не с вами? И
вообще как все? Жду с нетерпением письма. Что вы думаете предпринять дальше?
Только

\lp{20260604_194437600}

\noindent
уж, пожалуйста, не двигайтесь с места пока это рискованно. У нас жуть был
переполох из-за Heimbach. Ее должны были выслать, как и с всех австрийских и
германских подданных в северную губернию, но потом, для женщин это отменили. И
вообще у нас замечательно мило, предупредительно обращаются с немцами. Они все
сами заявляют об этом. Но с Heimbach дело осложнялось, потому что она не знает
русского языка и совершенно одна. Папаша и Коля переговорили с Мюллером и
Бокльманом, чтобы она могла присоединиться к этим семьям. Теперь она останется,
но ей надо искать себе работу, т.к. она не может оставаться у Юровских, которые
стали к ней дурно относиться.

\lp{20260604_194429422}

В данное время это очень трудно, и мы не знаем что из этого выйдет. Она тоже
страшно обрадовалась, что вы нашлись, и тоже напишет вам.

Очень тоскую по вас и хотела бы быть там же, где и вы. Неужели вы без Ядвиги?

Как вы доехали до Цюриха? Вероятно, стоя в проходе? Говорят, что это был
ужасный момент для переездов. Удалось ли вам взять свои вещи, а главное --
книги? Маня приехала совсем почти голая. Что вы там делаете? Мы, главным
образом, читаем газеты и делаем что-нибудь, имеющее отношение к событиям.
Почти все заняты этим же.

Обнимаю вас, мой любимый! Ваша Анюта.

Начиная с 16-го пишу вам ежедневно.
